\documentclass[12pt]{article}

\usepackage{bbm}
\usepackage{fancyhdr}
\usepackage{titling}
\usepackage{enumitem}
\usepackage{amsmath}
\usepackage{dsfont}
\usepackage{amsfonts}
\usepackage{amssymb}
\usepackage{amsthm}
\usepackage{geometry}
\usepackage{tikz}
\usepackage{algorithm}
\usepackage{algpseudocode}
\geometry{verbose,tmargin=1in,bmargin=1in,lmargin=1in,rmargin=1in}
\renewcommand{\vec}[1]{\mathbf{#1}}
\newcommand{\dec}{\text{dec}}
\newcommand{\divi}{\mathbin{\textnormal{div}}}
\title{Homework 1}
\author{Satvik Prasad}
\date{}
\newtheorem{lemma}{Lemma}[section]
\fancypagestyle{plain}{%
  \fancyhf{}
  \fancyhead[L]{\theauthor}
  \fancyhead[R]{\thedate}
}
\pagestyle{plain}
\makeatletter
\def\@maketitle{%
  \newpage
  \null
  \vskip 1em%
  \begin{center}%
    \let \footnote \thanks
    {\LARGE \@title \par}%
    \vskip 1em%
  \end{center}%
  \par
\vskip 1em}
\makeatother
\usepackage{cmbright}

\newcommand{\mat}[1]{\mathbf{#1}}

% ---------- Probability theory ----------
% Core operators (optional argument-free; use as \E{X}, \Var{X}, \Cov{X}{Y})
\DeclareMathOperator{\EOp}{\mathbb{E}}
\DeclareMathOperator{\VarOp}{Var}
\DeclareMathOperator{\CovOp}{Cov}
\DeclareMathOperator{\CorrOp}{Corr}

\renewcommand{\P}[1]{\mathbb{P}\left(#1\right)}        % \P{A}
\newcommand{\E}[1]{\EOp\left[#1\right]}                % \E{X}
\newcommand{\Var}[1]{\VarOp\left(#1\right)}            % \Var{X}
\newcommand{\Cov}[2]{\CovOp\left(#1,\,#2\right)}       % \Cov{X}{Y}
\newcommand{\Corr}[2]{\CorrOp\left(#1,\,#2\right)}     % \Corr{X}{Y}
\newcommand{\SD}[1]{\operatorname{SD}\left(#1\right)}  % \SD{X}

% Subscripted versions, e.g. \Esub{\theta}{X}, \Psub{\theta}{A}
\newcommand{\Esub}[2]{\EOp_{#1}\left[#2\right]}
\newcommand{\Psub}[2]{\mathbb{P}_{#1}\left(#2\right)}

% Conditioning, indicators, independence
\newcommand{\given}{\,\middle|\,}                      % \E{X \given Y}
\newcommand{\ind}[1]{\mathbbm{1}\left\{#1\right\}}     % \ind{X > 0}
\newcommand{\indep}{\perp\!\!\!\perp}                  % X \indep Y
\newcommand{\iid}{\overset{\text{iid}}{\sim}}          % X_i \iid \Normal(0,1)

% Convergence
\newcommand{\convp}{\xrightarrow{\;p\;}}               % in probability
\newcommand{\convd}{\xrightarrow{\;d\;}}               % in distribution
\newcommand{\convas}{\xrightarrow{\text{a.s.}}}        % almost surely
\newcommand{\convLp}[1]{\xrightarrow{L^{#1}}}          % in L^p

% Probability spaces
\newcommand{\F}{\mathcal{F}}                           % sigma-algebra
\newcommand{\B}{\mathcal{B}}                           % Borel sigma-algebra
\newcommand{\pspace}{(\Omega, \F, \mathbb{P})}

% Common distributions
\newcommand{\Normal}{\mathcal{N}}
\newcommand{\Bern}{\operatorname{Bernoulli}}
\newcommand{\Bin}{\operatorname{Binomial}}
\newcommand{\Pois}{\operatorname{Poisson}}
\newcommand{\Unif}{\operatorname{Uniform}}
\newcommand{\Exp}{\operatorname{Exp}}
\newcommand{\Geom}{\operatorname{Geometric}}
\newcommand{\Gam}{\operatorname{Gamma}}
\newcommand{\Beta}{\operatorname{Beta}}

% Number sets
\newcommand{\R}{\mathbb{R}}
\newcommand{\N}{\mathbb{N}}
\newcommand{\Z}{\mathbb{Z}}


\begin{document}
\subsection{Green Circle}
\begin{enumerate}
    \item[1] Let $n = 3k + r$ for some integer $k \geq 0,$ and integer $r \in \{0, 1, 2\}.$ Then, we know
    $$
    2^{n} - 1 \equiv 2^{3k + r} -1  \pmod 7.
    $$
    Factoring, we get 
    $$
    2^{n} -1 \equiv 8^{k} \cdot 2^{r} -1 \pmod 7 \equiv 2^{r} - 1 \pmod 7.
    $$
    Hence, this is only $\equiv 0 \pmod 7$ if and only if $r = 0,$ meaning $7 \mid 2^{n} - 1$ if and only if $n \equiv 0 \pmod 3.$ Thus, $n \in \{3, 6, 9, \ldots\}$ are the positive integers for which $7 \mid 2^{n} - 1.$

    Similarly, for $2^n + 1$ we get
    $$
    2^{n} + 1  \equiv 2^{r} + 1 \pmod 7,
    $$
    which is never $\equiv 0 \pmod 7$ for any $r \in \{0, 1, 2\}.$ Hence, $7$ never divides $2^{n} + 1$ for all positive integers $n.$

    \item[3] Define $n_0$ as the following:
    $$
    n_0 = \prod_{k = 1}^{10^{6} - 1} k.
    $$
    Then, the $10^6$ consecutive integers
    $$
    \{n_0, n_0 + 1, \ldots, n_0 + 10^{6} - 1\}
    $$
    are all composite, since the $k$-th entry (for $k > 1$) can be written as
    $$
    n_0 + k - 1 = \underbrace{(k-1)}_{> 0} \underbrace{\prod_{i \in \{1, 2, \ldots, 10^6\} \setminus\{k - 1\}} i}_{\in \mathbb N},
    $$
    and $n_0$ is trivially composite.

    Thus, there are a million consecutive composite integers.

    \item[4] Take $\gcd(3n + 1, 2n + 1).$ By the Euclidean algorithm, this is equivalent to $$
    G = \gcd(2n + 1, 3n + 1 \bmod 2n+ 1).
    $$
    Writing $3n + 1 = (2n + 1) + n,$ we can see $0 \leq n < 2n + 1,$ meaning $n$ is in residual range, and hence $3n + 1 \bmod 2n + 1 = n.$ Thus, we get
    $$
    G = \gcd(2n + 1, n) = 1
    $$
    as required. Therefore, $\gcd(2n + 1, 3n + 1) = 1,$ meaning $2n + 1, 3n + 1$ are coprime.
\end{enumerate}

\subsection{Blue Square}
\begin{enumerate}
    \item[7] Any such number $a$ can be written in the form
    $$
    a = 1 + 10^2 + \ldots + 10^{2k}
    $$
    for some integer $k \geq 0.$

    Thus, we can write
    $$
        a = \dfrac{(10^2)^{k} - 1}{10^{2} - 1} = \dfrac{(10^k - 1)(10^k + 1)}{99} = \dfrac{(1 + 10 + \ldots + 10^{k})(10^{k} + 1)}{11}.
    $$
    Since $k \geq 1$ (otherwise we have $a = 1$), we can write
    $$
    10^k + 1 = (11 - 1)^k + 1 = \sum_{r = 1}^{k} 11^{r} {k \choose r} = 11 \sum_{r = 0}^{k -1 } 11^{r}{k \choose r + 1},
    $$
    and hence
    $$
    a = \left (\sum_{r = 0}^{k} 10^r \right)\left(\sum_{r = 0}^{k-1} 11^r {k \choose r + 1}\right).
    $$
    If $k \geq 2,$ then both terms are $> 1,$ meaning $a$ is guaranteed to be composite. Hence, we must have $k =1,$ in which case $a = 101,$ which is a prime number.

    Thus, the only such prime is $101.$

    \item[8] Take cases on the parity of $n.$ If $n = 2k$ for some integer $k,$ we have
    $$
    n^4 + 4^n = 2^4 k^4 + 4^{2k} = 2(8k^4 + 2 \cdot 4^{2k - 1}).
    $$
    Since $n > 1,$ $k \geq 1,$ meaning the inner term is an integer, with $$
    8k^4 + 2 \cdot 4^{2k-1} \geq 8 + 2 \cdot 4^{2k-1} \geq 16,
    $$
    meaning we can write $n^4 + 4^n$ as the product of $2$ and some other factor $ > 1.$ Hence, $n^4 + 4^n$ is composite.

    Alternatively, $n = 2k + 1.$  Then, we have
    \begin{align*}
    n^4 + 4^n = (2k+1)^4 + 4^{2k + 1} &= (2k + 1)^4 + 2^{4k + 2} \\
    &= (2k+1)^4 + 4 (2^k)^4.
    \end{align*}

    We know 
    $$
    a^4 + 4b^4 = (a^2  + 2b^2)^2 -  4a^2b^2 = (a^2 + 2b^2 - 2ab)(a^2 + 2b^2 + 2ab),
    $$
    meaning we can write
    $$
    n^4 + 4^n = (a^2 + 2b^2 - 2ab)(a^2 + 2b^2 + 2ab)
    $$
    where $a, b$ are integers, and $a\geq 3,$ $b \geq 2.$ Then, 
    $$
    a^2 + 2b^2 + 2ab \geq 9 +32 + 24 = 65,
    $$
    and $$
    a^2 + 2b^2 - 2ab = (a+b)^2 + b^2 \geq 29.
    $$
    Hence, $n^4 + 4^n$ can be decomposed into a product of two factors of which neither are $1$, and hence $n^4 + 4n$ is composite in either case.
 
    \item[9] The Euler totient of $p^2$ is 
    $$
    \phi(p^2) = p^2 \left(1 - 1/p\right) = p(1-p),
    $$
    meaning for coprime $a, p^2$, we get
    $$
    a^{p^2 - p} \equiv 1 \pmod {p^2}.
    $$
    Equivalently, 
    $$
    a \cdot a^{p^2 - p - 1} \equiv 1\pmod {p^2},
    $$
    meaning if $\gcd(a, p) = 1,$ $a^{p^2 - p - 1}$ is the modular inverse of $a.$ Also, note that for $\gcd(a, p) = 1.$

    By the definition of the binomial coefficient,
    $$
    {2p \choose p} = \dfrac{2p(2p-1) \cdot \ldots \cdot  (p + 1)}{p \cdot \ldots \cdot 1} = \dfrac{2(2p - 1) \cdot \ldots \cdot (p + 1)}{(p- 1) \cdot \ldots \cdot 1}.
    $$
    Taking both sides mod $p^2,$ we have
    $$
        {2p \choose p} \equiv \dfrac{2(2p-1) \cdot \ldots \cdot (p + 1)}{(p-1) \cdot \ldots \cdot 1} \pmod {p^2}.
    $$

    Notice that if $a / b = c$ for integers $a, b, c$ then $a = bc.$ Under some modulo $n,$ we get 
    $$
    a \equiv bc \pmod n,
    $$
    and if $b^{-1}$ is the modular inverse of $b,$ then 
    $$
    ab^{-1} \equiv c \equiv a /b \pmod n.
    $$
    Hence, we know 
    \begin{align*}
    {2p \choose p} &\equiv 2(2p-1) \cdot  \ldots \cdot (p + 1) \cdot \prod_{k = 1}^{p-1} k^{p^2 - p - 1} \pmod {p^2} \\
    &\equiv 2 \prod_{k = 1}^{p-1} (k+p)k^{p^2 - p - 1} \pmod {p^2} \\
    &\equiv 2 \prod_{k = 1}^{p-1} k^{p^2 - p} + pk^{p^2 -p - 1} \pmod {p^2} \\
    &\equiv 2 \prod_{k = 1}^{p-1}1 + pk^{p^2 -p - 1} \pmod {p^2}.
    \end{align*}
    Since under modulo $p^2$, any terms that have a $p^2$ factor are congruent to $0$, we know
    \begin{align*}
    \kappa :=\prod_{k = 1}^{p - 1}(1 + pk^{-1}) \equiv \sum_{S \subseteq \{1, \ldots, p-1\}} \prod_{k \in S} pk^{-1} &\equiv \sum_{S \subseteq \{1, \ldots, p-1\}, |S| \leq  1} \prod_{k \in S} pk^{-1} \pmod{p^2} \\
    &\equiv 1 + \sum_{k = 1}^{p-1} pk^{-1} =  1 + p^{p-1} \sum_{k = 1}^{p -1} k. \pmod{p^2}
    \end{align*}
    When $p > 2,$ $p^2 \mid p^{p-1},$ and we trivially have $\kappa \equiv 1 \pmod {p^2}.$ Otherwise, when $p = 2$ we have $\kappa \equiv 3 \pmod {p^{2}}.$

    Thus, when $p > 2$ we know
    $$
    {2p \choose p} \equiv 2 \kappa \equiv 2 \pmod {p^2},
    $$
    and otherwise $$
    {2p \choose p } \equiv 2 \kappa \equiv 2 \cdot 3 \equiv 6   \equiv 2 \pmod{\underbrace{p^2}_{4}}.
    $$

    Therefore, in all cases we have
    $$
    {2p \choose p} \equiv 2 \pmod{p^2}
    $$
    for every prime $p;$
\end{enumerate}

\subsection{Black Diamond}
\begin{enumerate}
    \item[11] We prove that for any $n > 1,$ $n \nmid 2^{n} - 1.$ Assume for the sake of contradiction that there is some $n > 1$ with $n \mid 2^{n} - 1.$ Consider the smallest prime factor $p$ of $n.$ Then, we have
    $$
    2^{n} - 1 = 2^{pM} - 1.
    $$
    By Fermat's little theorem,
    $$
    2^n - 1 \equiv (2^{p})^M - 1 \equiv 2^M - 1 \pmod p.
    $$
    Additionally, since $p \mid n,$ and $n \mid 2^{n} -1,$ we know $2^{n} -1 \equiv 0 \pmod p,$ and thus $$
    2^{M} - 1 \equiv 0 \pmod {p} \Rightarrow 2^{M} \equiv 1 \pmod p.
    $$
    
    By Fermat's little theorem, we also know
    $$
    2^{p-1} \equiv 1 \pmod p.
    $$
    Therefore, $$2^{\gcd(M, p - 1)} \equiv 1 \pmod p,$$ since we can write $\gcd(M, p - 1)$ with Bezout coefficients. Since $p$ is the smallest prime factor of $M,$ no prime factor of $p - 1$ can be common with factors of $M,$ and hence $\gcd(M, p - 1) = 1.$ Hence, we get $2 \equiv 1 \pmod p,$ which is a contradiction for all prime $p.$

    Thus, $n \nmid 2^{n} -1 $ by contradiction.

    \item[12] Since $a - b \in \mathbb Z,$ we know $a = b + M$ for some integer $M \in \mathbb Z.$ Additionally, we have
    $$
    a^2 - b^2 = (a-b)(a + b) = M(a + b) \in \mathbb Z,
    $$
    or equivalently, 
    $$
    M(2b + M) = N
    $$
    for some integer $N,$ yielding $$
    b = \dfrac{N - M^2}{2M},\ a = \dfrac{N + M^2}{2M}
    $$
    for arbitrary integers $N, M \in \mathbb Z.$ Hence, $a, b$ are rational, meaning there exists a common denominator $d \in \mathbb Z$ such that
    $a = x/d,\ b = y/d$ for integers $x, y,$ with $\gcd(x, y,d) = 1.$

    Then, since $a^k - b^k \in \mathbb Z$ for all integers $k,$ we know
    $$
    d^k \mid x^k - y^k
    $$
    for all integers $k.$

    Consider some arbitrary prime factor $p \mid d.$ Then, we need the multiplicity of $p$ in $x^k - y^k$ to be $k$ for all integers $k.$ Factoring, we get
    $$
    x^k -y^k = (x-y)(x^{k - 1} + x^{k-2}y + \ldots + y^{k-1}).
    $$

    Since $x \equiv y,$ assume $x = y + p^s t$ for some integer $s$ and $p \nmid t.$ Then, $$x^k - y^k = (y + p^st)^k - y^k = \sum_{j = 1}^{k}{k \choose j}p^{sj}t^j y^{k -j}.$$
    When $k = p,$ we have
    \begin{align*}
    x^p - y^p = \sum_{j = 1}^{p}{p \choose j} p^{sj} t^j y^{p -j} &= p^{s} \sum_{j = 0}^{p - 1}{p \choose j + 1} p^{sj} t^{j + 1} y^{p-j - 1} \\
    &= p^{s+1} \sum_{j = 0}^{p-1} \dfrac{1}p{}{p \choose j + 1} p^{sj} t^{j + 1} y^{p -j - 1}.
    \end{align*}
    Since ${p \choose j + 1} / p \in \mathbb Z$ for all $0 \leq j \leq p - 1,$ clearly $p^{s + 1} \mid x^p - y^p.$ Additionally, we know
    $$
    x^p -y^p = p^{s + 1} \left (
    ty^{p-1} + \sum_{j = 1}^{p-1} \dfrac{1}{p}{p \choose j + 1} p^{sj} t^{j + 1} y^{p -j - 1}
    \right).
    $$
    Taking the second term modulo $p,$ we know
    $$
        ty^{p-1} + \sum_{j = 1}^{p-1} \dfrac{1}{p} {p \choose j + 1} p^{sj} t^{j + 1} y^{p -j - 1} \equiv ty^{p - 1}  \not \equiv 0 \pmod{p}
    $$
    since $p \nmid t$ and $\gcd(y, d) = 1.$ Thus, the multiplicity of $p$ in $x^p - y^p, $ $\kappa(x^p - y^p),$ is $s + 1,$ where $s = \kappa(x - y)$ is the multiplicity of $p$ in $x - y.$

    Recursively applying this rule with $x \to x^{p^{r - 1}}, y \to y^{p^{r - 1}}$, we know 
    $$
    x^{p^r} - y^{p^r} = (x^{p ^{r-1}})^{p} -(y^{p^{r}})^{p},
    $$
    meaning
    $$
    \kappa(x^{p^r} - y^{p^r}) = \kappa(x^{p^{r-1}} - y^{p^{r-1}}) + 1,
    $$
    and inductively,
    $$
\kappa(x^{p^r} - y^{p^r}) = \kappa(x - y) + r.
    $$

    Since for all $k = p^r,$ $d^{p^r} \mid x^{p^r} - y^{p^r},$ we know
    $$
    p^r \cdot \kappa(d) \leq \kappa(x^{p^r} - y^{p^r}) = \kappa(x - y) + r = s + r
    $$
    for some integer $s,$ and all integers $r.$ Rearranging, we get 
    $$
    r \leq \log_{p}\dfrac{s + r}{\kappa(d)},
    $$
    which cannot hold for all integers $r$, deriving a contradiction.

    Thus, no such $p \mid d$ can exist, i.e., $d = 1,$ meaning $a = x, b = y$ are integers.
\end{enumerate}

\subsection{Reflection}
For (12), it was incredibly tricky deriving some algebraic result for $a$, $b$ in terms of $N, M,$ so decided to prove that $a, b$ are rational instead. Then, since they're rational we can write them with a common denominator and $\gcd(\cdot, \cdot, \cdot) = 1,$ which seemed to suggest that the term $x^k - y^k$ (an integer) had to have factors of shared denominator $d$ scale exponentially, which could only occur if both $x$ and $y$ had that same denominator factor.

Nice to work with primes, so considered some prime factor of $d$ (notice that we RTP that $d = 1$, i.e, no such prime factor exists). Then, could prove that $d^k$ has prime factors scale exponentially, but $x^k - y^k$ only scales linearly under our assumption that neither $y$ nor $x$ contain prime factor $p.$
\end{document}